\begin{ledgroupsized}[r]{120mm}
\footnotesize
\pstart
\noindent\textbf{Überlieferung:}
\pend
\end{ledgroupsized}
\begin{ledgroupsized}[r]{114mm}
\footnotesize
\pstart \parindent -6mm
\makebox[6mm][l]{\textit{L}}%
Konzept:
LH~XXXV~9,~20 Bl.~2.
Ein Blatt 4\textsuperscript{o}.
Oberer, linker und rechter Rand beschnitten.
\href{\Scan{LH_35_09_20_002wz}}%
{Hälfte eines Wasserzeichens} 
am unteren Rand.
Zwei Seiten.
Auf 2~v\textsuperscript{o}, nach Ende des Konzepts (\refpassage{AE_1687_197-198_2a}{AE_1687_197-198_2b}), eigh.\ Rechnungen ohne erkennbaren Zusammenhang mit dem Text.
\pend
\end{ledgroupsized}
\begin{ledgroupsized}[r]{114mm}
\footnotesize
\pstart \parindent -6mm
\makebox[6mm][l]{\textit{l}}%
Reinschrift von \textit{L} (Hand von J.\,D.~Brandshagen%
\indexN{\textls{Brandshagen} (Brondhaguen), Jobst Dietrich 1659\textendash nach 1716}) %
mit Verbesserungen von Leibniz (\textit{Lil}):
LH~XXXV~9,~20 Bl.~3\textendash4.
Ein Bogen~2\textsuperscript{o}.
Ränder ausgefranst.
\href{\Scan{LH_35_09_20_003wz}}%
{Wasserzeichen} in der Mitte von Bl.~3.
Eineindrittel Seiten auf Bl.~3~r\textsuperscript{o} und 3~v\textsuperscript{o}.
Bl.~3~v\textsuperscript{o}\textendash4~r\textsuperscript{o} überliefern eine weitere Reinschrift durch Brandshagen: 
Leibnizens Rezension zu
\indexN{\textls{Spoleti}, Francesco, gest. 1712}\textsc{F.~Spoleti},
\cite{02095}\textit{De momento quo gravia nituntur deorsum per datum planum} \lbrack...\rbrack\ \textit{de secretione bilis in Hepate}, Venedig 1686. Die von Leibniz nicht veröffentlichte Rezension erscheint als N.\;\ref{41200} in diesem Band.
Bl.~4~v\textsuperscript{o} bis auf den eigh. Ordnungsvermerk „Mathematica“ leer.
\pend
\end{ledgroupsized}
\begin{ledgroupsized}[r]{114mm}
\footnotesize
\pstart
\parindent -6mm
\makebox[6mm][l]{\textit{E\textsuperscript{1}}}%
Erstdruck:
\cite{01023}\textit{AE} April 1687, S.\;197\,f.
(Unsere Druckvorlage.) Wie zu Leibnizens Zeiten üblich erschien die Rezension anonym.
\pend
\end{ledgroupsized}
\begin{ledgroupsized}[r]{114mm}
\footnotesize
\pstart
\parindent -6mm
\makebox[6mm][l]{\textit{E\textsuperscript{2}}}%
Abdruck (nach \textit{E\textsuperscript{1}}): 
\cite{01421}\textit{Opuscula omnia Actis eruditorum Lipsiensibus inserta, quae ad universam mathesin, physicam, medicinam, anatomiam, chirurgiam et philologiam pertinent}, Venedig 1740--1743,  Bd.~I, S.\;501\,f.
\pend%
\end{ledgroupsized}
\vspace{5mm}
\begin{ledgroup}
\footnotesize
\pstart
\noindent%
\textbf{Datierungsgründe:}
Die drei Zeugen der vorliegenden Rezension zu 
\indexN{\textls{Vanni}, Giovanni Francesco 1638\textendash1709}G.\,F.~Vannis anonym erschienenem Werk
\cite{02000}\textit{Exegeses physico-mathematicae, de momentis gravium, de vecte, ac de motu aequabiliter accelerato}, Rom 1685,
weisen eine lineare Entstehungsgeschichte auf.
(Siehe auch Leibnizens Besprechungen von Vannis bereits um 1684 erschienenem \cite{01995}\textit{Specimen} aus den \textit{Exegeses}: N.\;\ref{intro_vanni-specimen} und N.\;\ref{intro_vanni-demonstratio} in diesem Band.)
Der erste Textzeuge, das eigenhändige Konzept \textit{L}, 
diente dem Amanuensis
\indexN{\textls{Brandshagen} (Brondhaguen), Jobst Dietrich 1659\textendash nach 1716}J.\,D.~Brandshagen als unmittelbare Vorlage für seine Reinschrift \textit{l}, wie aus verschiedenen Fehlern in der Abschrift (Haplographien, Zeilensprüngen) klar hervorgeht.
Die Reinschrift, einschließlich der Verbesserungen durch Leibniz (Textschicht \textit{LiL}) war wiederum Vorlage für den Druck \textit{E\textsuperscript{1}}; weitere Textzeugen (etwa eine weitere Reinschrift von \textit{Lil} als \textit{unmittelbare} Vorlage für \textit{E\textsuperscript{1}}) sind nicht bekannt.
\pend
\pstart
Leibnizens Briefwechsel mit O.~Mencke\indexN{\textls{Mencke} (Menken, Menkenius, Menque), Otto 1644\textendash1707} %
gibt Auskunft über die Entstehung der Stücke zwischen Ende Oktober 1686 und Mitte Februar 1687.
Am \cite{02013}14. (24.) Oktober 1686 bestätigt Mencke die Zusendung einiger Bücher, 
darunter Vannis \cite{02000}\textit{Exegeses}, und bittet Leibniz um eine \glqq relation ad Acta\grqq\ des Werks (\textit{LSB} I, 4 N.\;499, hier S.\;601). 
Das Schreiben liefert daher einen Terminus post quem für den Textzeugen \textit{L} und insgesamt für N.\;\ref{41199}.
In seinem \cite{02014}Brief vom 2. (12.) März 1687 (\cite{00115}\textit{LSB} I, 4 N.\;521)
bedankt sich Mencke\indexN{\textls{Mencke} (Menken, Menkenius, Menque), Otto 1644\textendash1707}
bei Leibniz für dessen (nicht erhaltenes) Schreiben vom 11. (21.) Februar und für die \glqq beygefügten recensionen\grqq, die im Aprilheft der \cite{01023}\textit{Acta Eruditorum} erscheinen sollen. 
Die \textit{Acta} brachten tatsächlich Leibnizens Besprechungen von
\indexN{\textls{Gregory} (Gregorie), David 1625\textendash1720}%
D.~Gregory, \textit{Exercitatio geometrica de dimensione figurarum}, Edinburgh 1684 (dort S.\;196\,f.)
und von \indexN{\textls{Vanni}, Giovanni Francesco 1638\textendash1709}Vannis 
\cite{02000}\textit{Exegeses} (dort S.\;197\,f., unser Textzeuge \textit{E\textsuperscript{1}}).
Die verbesserte Reinschrift \textit{Lil} von N.\;\ref{41199} (bzw. eine nicht bekannte Abschrift davon)
zählte aller Wahrscheinlichkeit nach zu den am 21. Februar verschickten "recensionen". Damit liegt ein Terminus ante quem für die Abfassung der Reinschrift vor.
\pend
\pstart
Die eigenhändigen Rechnungen am Ende des Textzeugen \textit{L} (Bl.~2~v\textsuperscript{o}, hier \refpassage{AE_1687_197-198_2a}{AE_1687_197-198_2b}) 
weisen zwar keinen erkennbaren Zusammenhang mit dem Text auf;
die zweite und dritte Spalte lassen sich jedoch als Hilfsrechnungen zum in einer benachbarten Signatur (LH~XXXV~9,~21 Bl.~1\textendash2) überlieferten Konzept \cite{01383}\textit{Leges Concursuum} identifizieren (\textit{LSB}~VIII,~3 N.\;67). 
Daher tragen die Rechnungen, dank der im Briefwechsel gut dokumentierten Entstehung von N.\;\ref{41199},
zur Datierung der \textit{Leges Concursuum} bei
(siehe die Datierungsbegründung in \textit{LSB}~VIII,~3 S.\;735 und die Hilfsrechnungen in der Marginalie zu S.\;742 Z.~2\textendash5.)
\pend 
\end{ledgroup}
\selectlanguage{latin}
\vspace{8mm}
\pstart
\centering
\edlabel{AE_1687_197-198_1a}%
\edtext{}{%
\xxref{AE_1687_197-198_1a}{AE_1687_197-198_1b}%
\lemma{}%
\Bfootnote{%
\href{\Scan{LH_35_09_20_002r}}%
{\lbrack2~r\textsuperscript{o}\rbrack} \textit{Exegeses physico-Mathematicae} \lbrack...\rbrack\ Romae ex typogr. Rev. Cam. Apost. 1685. 8\textsuperscript{o}.~\textit{L}
\textbullet~%
\href{\Scan{LH_35_09_20_003r}}%
{\lbrack3~r\textsuperscript{o}\rbrack} \textit{Exegeses physico-Mathematicae} \lbrack...\rbrack\ \textit{de motu aequabiliter accelerato}. Romae ex typogr. Rev. Cam. Apost. 1686. 8\textsuperscript{o}.~\textit{l}
\textbullet~%
\lbrack S.\;197\rbrack~\textit{EXEGESES} \lbrack...\rbrack\ Apost. \bars{1686 \textit{ändert Hrsg.}}, in 8.~\textit{E\textsuperscript{1}}}}%
\lbrack S.\;197\rbrack
\hfill
\edtext{%
\cite{02000}\textit{EXEGESES PHYSICO-MATHEMATICAE}
\hfill
\phantom{\lbrack S.\;197\rbrack}
\linebreak
\textit{de Momentis gravium, de Vecte, de Motu aequabiliter accelerato}.\linebreak
Romae, ex typograph. Rev. Cam. Apost.
\lbrack1685\rbrack,}{%
\lemma{\textit{EXEGESES} \lbrack...\rbrack\ \lbrack1685\rbrack}%
\Cfootnote{%
\indexN{\textls{Vanni}, Giovanni Francesco 1638\textendash1709}\textsc{G.\,F.~Vanni}, 
\cite{02000}\textit{Exegeses physico-mathematicae, de momentis gravium, de vecte, ac de motu aequabiliter accelerato}, Rom 1685. Das Buch erschien anonym.
Das Erscheinungsdatum wird in \textit{L} richtig als 1685 angegeben, in \textit{l} und \textit{E\textsuperscript{1}} hingegen als 1686.}}
in 8.\edlabel{AE_1687_197-198_1b}
\pend%
\vspace{0.5\baselineskip}%
\pstart  
\noindent
Quemadmodum laudandi
\edtext{sunt,}{%
\lemma{sunt}%
\Bfootnote{\textit{L~u.~l}
\repeatlemma,~\textit{E\textsuperscript{1}}}}
qui locupletant
\edtext{scientias:}{%
\lemma{scientias}%
\Bfootnote{\textit{L}
\repeatlemma,~\textit{l}
\repeatlemma:~\textit{E\textsuperscript{1}}}} 
ita illorum quoque non spernenda est industria,
qui
\edtext{saltem}{%
\lemma{}%
\Bfootnote{%
saltem %
\textit{erg.~L}}}
elegantia
\edtext{quaedam}{%
\lemma{quaedam}%
\Bfootnote{,~\textit{l}
\repeatlemma~\textit{L~u.~E\textsuperscript{1}}}}
dubia
\edtext{proponunt}{%
\lemma{proponunt}%
\Bfootnote{,~\textit{L}
\repeatlemma~\textit{l}}}
contra sententias
\edtext{receptas.\indexS{sententia recepta} Nam}{%
\lemma{receptas}%
\Bfootnote{%
,~%
\textit{(1)}~etiamsi ea essent verae et demonstratae %
\textit{(2)}~licet %
\textit{(3)}~sive %
\textit{(4)}~nam~\textit{L}
\repeatlemma; nam~\textit{l}
\repeatlemma. Nam~\textit{E\textsuperscript{1}}}}
si falsae sunt 
\edtext{istae,}{%
\lemma{}%
\Bfootnote{%
istae, %
\textit{erg.~L}}}
veritas
\edtext{detegitur. Sin}{%
\lemma{detegitur}%
\Bfootnote{%
,~%
\textit{(1)}~si %
\textit{(2)}~sin~\textit{L}
\repeatlemma; Sin~\textit{l}
\repeatlemma. Sin~\textit{E\textsuperscript{1}}}}
verae et demonstratae 
\edtext{sint,}{%
\lemma{}%
\Bfootnote{%
sint \textit{erg.~L}
\repeatlemma,~\textit{l}}}
saltem ex solutione dubii nova quaedam
\edtext{lux effulget.}{%
\lemma{}%
\Bfootnote{%
lux \bars{emergit, et~\textit{gestr.}} %
effulget~\textit{L}}}
\pend
\begin{edfigure}%
\includegraphics[scale=1]{%
gesamttex/edit_VIII-4/images/AE_1687_197-198_d1_002r.pdf}
\edcaption{\lbrack\textit{Fig.~1, gestr.; L (Bl.~2~r\textsuperscript{o})}\rbrack}%
\end{edfigure}
\pstart
\edlabel{AE_1687_197-198_5a}%
\edtext{}{%
\xxref{AE_1687_197-198_5a}{AE_1687_197-198_5b}%
\lemma{Primus \lbrack...\rbrack\ aequalia}%
\Cfootnote{%
\indexN{\textls{Tartaglia} (Tartalea), Niccolò 1499 o. 1500\textendash1557}%
\textsc{N.~Tartaglia}, %
\cite{02016}\textit{Quesiti et inventioni diverse}, %
2.~Ausgabe, Venedig 1554. %
Libro Ottavo, Quesito XLII Propositione XV, %
Bl.~97~r\textsuperscript{o}\textendash97~v\textsuperscript{o}.}}%
Primus Nicolaus
\edtext{\indexN{\textls{Tartaglia} (Tartalea), Niccolò 1499 o. 1500\textendash1557}Tartalea,}{%
\lemma{Tartalea}%
\Bfootnote{\textit{l}
\repeatlemma,~\textit{L~u.~E\textsuperscript{1}}}}
quantum constat,
\edtext{invenit gravium momenta\indexS{momentum gravis in plano inclinato}}{%
\lemma{invenit}%
\Bfootnote{%
\textit{(1)}~Gravia in %
\bars{ejusdem elevationis posit\textlangle a\textrangle~\textit{erg.}} %
planis inclinatis habere momenta reciproce ut plana %
\bars{sive ut lineas~\textit{erg.}} %
, seu ut ipse loquitur, duorum gravium aequalia esse momenta si %
\textit{(2)}~Gravia %
\textit{(3)}~Gravium momenta~\textit{L}
\repeatlemma
\textit{(1)}~aequalia %
\textit{(2)}~Gravium momenta~\textit{l}
}}
\edtext{in planis,}{%
\lemma{in}%
\Bfootnote{
\textit{(1)}~plan %
\textit{(2)}~lineis rectis %
\textit{(3)}~planis~\textit{L}
\repeatlemma planis~\textit{l}
\repeatlemma planis,~\textit{E\textsuperscript{1}}}}
sive potius lineis rectis\indexS{linea recta inclinata}
\edtext{inclinatis,}{%
\lemma{inclinatis}%
\Bfootnote{\textit{L~u.~l}
\repeatlemma,~\textit{E\textsuperscript{1}}}} 
esse reciproce
ut ipsas lineas
\edtext{inclinatas,\indexS{linea recta inclinata}}{%
\lemma{inclinatas}%
\Bfootnote{\textit{L}
\repeatlemma,~\textit{l}}}
positis
\edtext{gravibus\indexS{grave} aeque elevatis,\indexS{gravia aeque elevata}}{%
\lemma{gravibus}%
\Bfootnote{%
\textit{(1)}~ejusdem %
\textit{(a)}~inclina %
\textit{(b)}~eleva %
\textit{(2)}~aeque elevatis;~\textit{L}
\repeatlemma aeque elevatis,~\textit{l}}}
vel ut
\edtext{\indexN{\textls{Tartaglia} (Tartalea), Niccolò 1499 o. 1500\textendash1557}Tartalea rem}{%
\lemma{Tartalea}%
\Bfootnote{%
\textit{(1)}~loquitur; si %
\textit{(2)}~rem~\textit{L}}}
eleganter
\edtext{exprimit: si gravia aeque elevata,\indexS{gravia aeque elevata} 
inter se sint ut ipsae rectae inclinatae,\indexS{linea recta inclinata} 
momenta\indexS{momentum gravis} eorum fore aequalia.%
\edlabel{AE_1687_197-198_5b}}{%
\lemma{exprimit}%
\Bfootnote{%
\textit{(1)}~gravium %
\textit{(a)}~aeque le %
\textit{(b)}~momenta esse ut ipsas rectas inclinat
\textit{(2)}~gravium aeque elevatorum momenta esse ut ipsas rectas inclinatas. %
\textit{(3)}~si %
\textit{(4)}~. Si gravia aeque elevata %
\textit{(a)}~ess %
\textit{(b)}~sint
\textit{(c)}~inter se \lbrack...\rbrack\ fore aequalia.~\textit{L}
\repeatlemma. Si gravia aeque elevata inter \lbrack...\rbrack\ %
inclinatae momenta eorum fore aequalia.~\textit{l}
\repeatlemma: \lbrack...\rbrack\ aequalia.~\textit{E\textsuperscript{1}}}}
Hoc postea
\edtext{\indexN{\textls{Stevin} (Stevinus), Simon 1548\textendash1620}Stevinus,}{%
\lemma{Stevinus}%
\Cfootnote{%
\indexN{\textls{Stevin} (Stevinus), Simon 1548\textendash1620}%
\textsc{S.~Stevin}, %
\cite{00196}\textit{De beghinselen der weeghconst}, %
Leiden 1586. %
Buch~1, Teil~2, Theorem~XI Prop.~XIX mit Koroll.~I, bes.\ S.\;40\textendash42
(\cite{02084}\textit{SPW} I, S.\;174\textendash179).}} 
\edtext{\indexN{\textls{Galilei} (Galilaeus, Galileus), Galileo 1564\textendash1642}Galilaeus,}{%
\lemma{Galilaeus}%
\Bfootnote{\textit{L~u.~l}
\repeatlemma,~\textit{E\textsuperscript{1}}%
}%
\Cfootnote{%
\indexN{\textls{Galilei} (Galilaeus, Galileus), Galileo 1564\textendash1642}%
\textsc{G.~Galilei}, %
\cite{00050}\textit{Discorsi}, %
Giornata Terza, De motu naturaliter accelerato; \glqq Aggiunta postuma dell'Autore\grqq\ %
nach dem Scholion zum 2.~Theorem. Der Zusatz erschien erst in der posthumen Werkausgabe: \cite{01084}\textit{Opere}, Bologna 1656, Bd.~II, S.\;132\lbrack ter\rbrack\textendash134 (\cite{00048}\textit{GO} VIII, S.\;216\textendash219 im Apparat).}}
\edtext{aliique excellentes viri}{%
\lemma{aliique excellentes viri}%
\Cfootnote{\indexN{\textls{Vanni}, Giovanni Francesco 1638\textendash1709}%
Vanni kritisiert neben \indexN{\textls{Stevin} (Stevinus), Simon 1548\textendash1620}Stevin 
und \indexN{\textls{Galilei} (Galilaeus, Galileus), Galileo 1564\textendash1642}Galilei 
auch Mersenne,\indexN{\textls{Mersenne} (Mersennus), Marin 1588\textendash1648} %
Dechales\indexN{\textls{Dechales} (Chalesius), Claude François Milliet 1621\textendash1678} %
und Borelli.\indexN{\textls{Borelli} (Borellus), Giovanni Alfonso 1608\textendash1679} %
Siehe \cite{02000}\textit{Exegeses}, pars I, S.\;4\,f.}}
prosecuti
\edtext{sunt, nonnulli}{%
\lemma{sunt,}%
\Bfootnote{%
\textit{(1)}~et compluribus demonstrationibus muniverunt. Nihilominus doctissimus autor haec %
\textit{(2)}~nonnulli~\textit{L}}}
et demonstrationibus
\edtext{varie}{%
\lemma{}%
\Bfootnote{varie \textit{erg.~L}}}
\edtext{muniverunt.}{%
\lemma{muniverunt}%
\Bfootnote{%
;~\textit{l}
\repeatlemma.~\textit{L~u.~E\textsuperscript{1}}}}
At doctissimo \indexN{\textls{Vanni}, Giovanni Francesco 1638\textendash1709}autori 
harum \cite{02000}\textit{Exegesium} in mentem venit 
\edlabel{AE_1687_197-198_6a}%
\edtext{}{%
\xxref{AE_1687_197-198_6a}{AE_1687_197-198_6b}%
\lemma{difficultas \lbrack...\rbrack\ proposuit}%
\Cfootnote{%
\indexN{\textls{Vanni}, Giovanni Francesco 1638\textendash1709}%
\textsc{G.\,F. Vanni}, \cite{01995}\textit{Specimen libri de momentis gravium}, in \cite{01023}\textit{AE} November 1684, S.\;511\textendash513. Der Aufsatz erschien anonym und war von einer ebenfalls anonymen Leibniz'schen \glqq censura\grqq\ begleitet, siehe N.\;\ref{55606} in diesem Band. In \textit{L} und \textit{l} dienen die Punkte als Platzhalter für die Seitenangabe.}}%
difficultas perelegans, quam in \cite{01995}\textit{Specimine} aliquo jam his 
\edtext{\cite{01023}\textit{Actis} a.\;1684. m.\,Nov. pag.\;511. inserto}{%
\lemma{\textit{Actis}}%
\Bfootnote{%
.~.~.~.~.~.~.~. inserto \textit{L~u.~l}
\repeatlemma a.\;1684. m.\,Nov. pag.\;511. inserto~\textit{E\textsuperscript{1}}}}
\edtext{proposuit,\edlabel{AE_1687_197-198_6b}}{%
\lemma{proposuit}%
\Bfootnote{\textit{l}
\repeatlemma,~\textit{L~u.~E\textsuperscript{1}}}}
\edtext{cui responderunt aliquot viri in his %
\edtext{studiis versati.}{%
\lemma{studiis}%
\Bfootnote{%
\textit{(1)}~insignes %
\textit{(2)}~versati.~\textit{L}}}%
}{%
\lemma{cui \lbrack...\rbrack\ versati}%
\applabel{FnC:AE_1687_197-198_demonstratio-geometrica}%
\Cfootnote{In den \cite{01023}\textit{Acta Eruditorum} erschienen folgende Repliken auf 
\indexN{\textls{Vanni}, Giovanni Francesco 1638\textendash1709}%
Vannis \cite{01995}\textit{Specimen}: %
\indexN{\textls{Kochański}, Adam Adamandus 1631\textendash1700}\textsc{A.\,A. Kochański}, %
\cite{01996}\textit{Consideratio speciminis libri de momentis gravium, Authore J.\,F.\,V. Lucensi}, %
Juni 1685, S.\;262\textendash265;
\indexN{\textls{Leibniz} (Leibnitius, GGL), Gottfried Wilhelm 1646\textendash1716}\textsc{Leibniz}, \textit{Demonstratio geometrica regulae apud Staticos receptae}, %
November 1685, S.\;501\textendash505 (N.\;\ref{49409} in diesem Band);
\indexN{\textls{Bernoulli}, Jacob 1655\textendash1705}\textsc{Jacob Bernoulli}, %
\cite{02003}\textit{Solutio difficultatis contra propositionem quandam mechanicam Autore I.\,F.\,V. Lucensi propositae}, 
Februar 1686, S.\;96\textendash98.}}
Nos judicium nobis non 
\edlabel{AE_1687_197-198_3a}%
\edtext{}{%
\xxref{AE_1687_197-198_3a}{AE_1687_197-198_3b}%
\lemma{arrogamus}%
\Bfootnote{%
, 
\textit{(1)}~illud tamen
\textit{(a)}~ex
\textit{(b)}~fortasse cum aliis possemus non injuria expectare,
\textit{(aa)}~non t
\textit{(bb)}~si sententiam receptam \bars{adhuc~\textit{erg.}} infirmam esse arbitratur
\bars{vir egregius~\textit{erg.}} non tantum
\textit{(aaa)}~sa
\textit{(bbb)}~etiam 
\textit{(aaaa)}~pro 
\textit{(bbbb)}~adductis pro ea demonstrationibus satisfacere velit
\textit{(aaaaa)}~; \bars{ex \textit{verseh.\ erhalten}} his ultra inprimis
\textit{(bbbbb)}~. Et
\textit{(aaaaa-a)}~quidem
\textit{(bbbbb-b)}~ex his
\textit{(ccccc-c)}~ex illis facilis esse videtur, quae pro
\textit{(ccccc)}~et stringere videtur, 
quae proposita est in \textit{Actis} nostris .~.~.~.~.~.~.~.
\textit{(ddddd)}~Et quidem ex illis, qui in \textit{Actis} nostris
\textit{(aaaaa-a)}~difficultati \textit{Speciminis}
\textit{(bbbbb-b)}~\textit{Specimini} respondere sunt conati unus attulit demonstrationem aliquam novam quae hoc habet peculiare, quod tota Geometrica est,
\textit{(aaaaa-aa)}~tantu
\textit{(bbbbb-bb)}~nisi quod ex Mechanicis assumit 
\textit{(ccccc-cc)}~ex Mechanicis tantum supponendo
\textit{(aaaaa-aaa)}~quae Archimedes habet de
\textit{(bbbbb-bbb)}~conc
\textit{(ccccc-ccc)}~centrum gravitatis
\textit{(ddddd-ddd)}~corpor
\textit{(eeeee-eee)}~regulam aequilibrii posita
\textit{(fffff-fff)}~corpo 
\textit{(ggggg-ggg)}~notissima illa
\textit{(hhhhh-hhh)}~notissimum illud quod corpora sunt
\textit{(iiiii-iii)}~notissima illa \bars{principia~\textit{erg.}}, quod corpora
\textit{(aaaaa-aaaa)}~descendere
\textit{(bbbbb-bbbb)}~sunt in aequilibrio, quando centrum eorum descendere non potest; et quando reciproce sunt ut distantiae a sustentaculo.
\textit{(2)}~sed integrum lectori relinquimus. Caeterum \bars{clarissimus \textit{gestr.}} autor
\textit{(a)}~in his \textit{Exegesibus}
\textit{(b)}~in harum \textit{Exegesium} prima~\textit{L}
\repeatlemma \lbrack...\rbrack\ prima~\textit{l}%
}%
\lemma{arrogamus \lbrack...\rbrack\ prima}%
\Cfootnote{%
Leibnizens Anspielung \glqq demonstrationem aliquam novam\grqq\ in der Variante \textit{(ddddd)} gilt seiner \textit{Demonstratio Geometrica}; vgl. die Erläuterung zu \refnote{FnC:AE_1687_197-198_demonstratio-geometrica}.
Zu \indexN{\textls{Archimedes} 287\textendash212 v. Chr.}Archimedes siehe
\cite{01010}\textit{De planorum aequilibris libri duo}, lib.~I, prop.~VI.}}%
arrogamus, sed integrum lectori relinquimus.
\pend
\pstart
Caeterum Autor in harum \cite{02000}\textit{Exegesium} prima%
\edlabel{AE_1687_197-198_3b}
non tantum receptam sententiam\indexS{sententia recepta}
impugnat, sed et novam stabilire
\edtext{tentat:}{%
\lemma{tentat}%
\Bfootnote{%
,~\textit{L}
\repeatlemma;~\textit{l}
\repeatlemma~:~\textit{E\textsuperscript{1}}}}
\edlabel{AE_1687_197-198_8a}%
\edtext{}{%
\xxref{AE_1687_197-198_8a}{AE_1687_197-198_8b}%
\lemma{putat \lbrack...\rbrack\ horizontale}%
\Cfootnote{%
\indexN{\textls{Vanni}, Giovanni Francesco 1638\textendash1709}%
\textsc{Vanni}, 
\cite{02000}\textit{Exegeses}, pars I, Prop.~XVI, S.\;25.}}%
putat scilicet
\edtext{momenta gravium in planis declivibus\indexS{momentum gravis in plano declivi}}{%
\lemma{momenta}%
\Bfootnote{%
\textit{(1)}~planorum %
\textit{(2)}~gravium in planis %
\textit{(a)}~inclinatis %
\textit{(b)}~declivibus~\textit{L}}}
esse ut
\edtext{angulos,\indexS{angulus plani ad planum horizontale} quos plana\indexS{planum declive} faciunt ad planum horizontale;\indexS{planum horizontale}\edlabel{AE_1687_197-198_8b}}{%
\lemma{angulos}%
\Bfootnote{%
\textit{(1)}~elevationis; momenta autem quibus pre %
\textit{(2)}~quos plana faciunt ad %
\textit{(a)}~horizontem, %
\textit{(b)}~planum horizontale,~\textit{L}
\repeatlemma \lbrack...\rbrack\ horizontale;~\textit{l}
\repeatlemma, \lbrack...\rbrack\ horizontale;~\textit{E\textsuperscript{1}}}}
\edlabel{AE_1687_197-198_7a}%
\edtext{}{%
\xxref{AE_1687_197-198_7a}{AE_1687_197-198_7b}%
\lemma{momenta \lbrack...\rbrack\ evanescant}%
\Cfootnote{%
\cite{02000}a.\,a.\,O., Prop.~XV, S.\;23\textendash25.}}%
momenta\indexS{momentum quo planum declive premitur} autem quibus ipsa 
\edtext{plana\indexS{planum declive} premuntur,}{%
\lemma{plana}%
\Bfootnote{%
\textit{(1)}~prementur %
\textit{(2)}~premuntur~\textit{L}
\repeatlemma premuntur,~\textit{l}}}
esse ut horum angulorum\indexS{complementum anguli elevationis}
\edtext{complementa: idque}{%
\lemma{complementa}%
\Bfootnote{%
. %
\textit{(1)}~quod %
\textit{(2)}~idque~\textit{L}
\repeatlemma; idque~\textit{l}
\repeatlemma: idque~\textit{E\textsuperscript{1}}}}
\edtext{pag.~24. probare instituit ex eo, quod plano}{%
\lemma{pag.}%
\Bfootnote{24.
\textit{(1)}~probat %
\textit{(2)}~probare %
\textit{(a)}~instituit %
\textit{(b)}~instituit ex eo quod %
\textit{(aa)}~quo magis %
\textit{(bb)}~plano~\textit{L}
\repeatlemma \lbrack...\rbrack\ eo quod plano:~\textit{l}
\repeatlemma \lbrack...\rbrack\ eo, quod plano~\textit{E\textsuperscript{1}}}}
prorsus
\edtext{jacente\indexS{planum prorsus jacens} totalis}{%
\lemma{jacente}%
\Bfootnote{%
totalis~\textit{L}
\repeatlemma, 
\textit{(1)}~talis~\textit{l} %
\textit{(2)}~totalis~\textit{Lil}
\repeatlemma totalis~\textit{E\textsuperscript{1}}}}
est
\edtext{pressio,\indexS{pressio totalis globi in planum} et erecto\indexS{planum magis magisque erectum}}{%
\lemma{pressio}%
\Bfootnote{
,~%
\textit{(1)}~quae erecto plan %
\textit{(2)}~crescente %
\textit{(3)}~et erecto~\textit{L}
\repeatlemma; et erecto~\textit{l}
\repeatlemma, et erecto~\textit{E\textsuperscript{1}}}}
magis magisque
\edtext{plano,}{%
\lemma{plano}%
\Bfootnote{\textit{L~u.~l}
\repeatlemma,~\textit{E\textsuperscript{1}}}}
\edtext{seu decrescente angulo complementi elevationis,\indexS{angulus complementi elevationis}}{%
\lemma{seu}%
\Bfootnote{%
\textit{(1)}~angulo complementi elevationis decrescente %
\textit{(2)}~decrescente angulo complementi elevationis,~\textit{L}}}
continue
\edtext{pres\lbrack S.\;198\rbrack sio\indexS{pressio globi in planum}}{%
\lemma{}%
\Bfootnote{pressio \textit{erg.~L}}}
imminuitur, donec plano prorsus perpendiculariter erecto,\indexS{planum prorsus perpendiculariter erectum} simul et
angulus complementalis\indexS{angulus complementalis elevationis} et pressio\indexS{pressio globi in planum} 
\edtext{globi}{%
\lemma{}%
\Bfootnote{globi \textit{erg.~L}}}
in 
\edtext{planum, evanescant.\edlabel{AE_1687_197-198_7b} Licet}{%
\lemma{planum,}%
\Bfootnote{%
\textit{(1)}~evanescat %
\textit{(2)}~evanescant. \href{\Scan{LH_35_09_20_002v}}%
{\lbrack2~v\textsuperscript{o}\rbrack} %
\textit{(a)}~Sed cum %
\textit{(b)}~Licet~\textit{L}}}
autem non semper 
\edtext{alias}{%
\lemma{}%
\Bfootnote{alias \textit{erg.~L}}}
\edtext{sequatur,}{%
\lemma{sequatur}%
\Bfootnote{\textit{L~u.~l}
\repeatlemma,~\textit{E\textsuperscript{1}}}}
ea quae simul incipiunt, crescunt
\edtext{et desinunt,}{%
\lemma{et}%
\Bfootnote{%
desinunt~\textit{L}
\repeatlemma %
\textit{(1)}~desinant~\textit{l} %
\textit{(2)}~desinunt~\textit{Lil}
\repeatlemma desinunt,~\textit{E\textsuperscript{1}}}}
esse proportionalia, ut multis exemplis
\edtext{constat; eadem enim ratio et ad sinus angulorum applicari potuisset: videtur tamen aestimatio per angulos\indexS{aestimatio momenti per angulos} 
\indexN{\textls{Vanni}, Giovanni Francesco 1638\textendash1709}%
Autori placuisse,}{%
\lemma{constat}%
\Bfootnote{%
\bars{; eadem \lbrack...\rbrack\ potuisset~\textit{erg.}}. %
\textit{(1)}~Videt %
\textit{(a)}~per %
\textit{(b)}~aestimatio \bars{per~\textit{erg.}}  angulos %
\textit{(2)}~Videtur tamen aestimatio per angulos autori placuisse~\textit{L}
\repeatlemma; \lbrack...\rbrack\ potuisset; Videtur \lbrack...\rbrack\ autori placuisse~\textit{l}
\repeatlemma; \lbrack...\rbrack\ placuisse,~\textit{E\textsuperscript{1}}}}
\edtext{quod concinne}{%
\lemma{quod}%
\Bfootnote{%
\textit{(1)}~commode %
\textit{(2)}~concinne~\textit{L}}}
admodum
\edtext{schemati in \cite{01023}\textit{Actis} nostris
\edtext{l.\,c.}{%
\lemma{}%
\Bfootnote{l.\,c.~\textit{fehlt} \textit{L~u.~l}}}
exhibito,}{%
\lemma{schemati \lbrack...\rbrack\ exhibito}%
\Cfootnote{%
\indexN{\textls{Vanni}, Giovanni Francesco 1638\textendash1709}%
\textsc{Vanni}, \cite{01995}\textit{Specimen} (N.\;\ref{55606} in diesem Band).}}
\edtext{quo globus}{%
\lemma{quo}%
\Bfootnote{%
\textit{(1)}~duo plana %
\textit{(2)}~globus~\textit{L}}}
duobus planis\indexS{planum} angulum rectum facientibus incumbit,\indexS{globus duobus planis incumbens} satisfacere 
\edtext{videatur. Licet alioqui geometrae\indexS{geometra} in}{%
\lemma{videatur}%
\Bfootnote{%
\textit{(1)}~. Licet alioqui G %
\textit{(2)}~, licet alioqui Geometrae \bars{omnes~\textit{gestr.}} in~\textit{L}}}
applicatione\indexS{applicatio linearum ad naturam rerum}
\edtext{linearum}{%
\lemma{}%
\Bfootnote{linearum \textit{erg.~L}}}
ad naturam
\edtext{rerum}{%
\lemma{}%
\Bfootnote{rerum \textit{erg.~L}}}
\edtext{instituenda,}{%
\lemma{instituenda}%
\Bfootnote{\textit{fehlt}~\textit{L~u.~l}, \textit{erg.~Lil}
\repeatlemma,~\textit{E\textsuperscript{1}}}}
\edtext{potius sinuum quam angulorum}{%
\lemma{potius}%
\Bfootnote{%
\textit{(1)}~sinus quam angulos %
\textit{(2)}~sinuum quam angulorum~\textit{L}}} 
\edtext{proportionem cum successu observare soleant. Si}{%
\lemma{proportionem}%
\Bfootnote{%
\bars{%
\textit{(1)}~feliciter %
\textit{(2)}~cum successu~\textit{erg.}}
observare soleant, \bars{ut exemplo refractionum patet~\textit{gestr.}}. Si~\textit{L}}}
vero duo
\edtext{plana,\indexS{plana quibus globus incumbet}}{%
\lemma{plana}%
\Bfootnote{\textit{L~u.~l}
\repeatlemma,~\textit{E\textsuperscript{1}}}}
quibus globus 
\edtext{incumberet, angulum}{%
\lemma{incumberet}%
\Bfootnote{%
\textit{(1)}~a globo premantur ex his %
\textit{(a)}~autoribus %
\textit{(b)}~autoris %
\textit{(2)}~angulum~\textit{L}
\repeatlemma, angulum~\textit{l}}}
inter se facerent
\edtext{non rectum, \indexN{\textls{Vanni}, Giovanni Francesco 1638\textendash1709}%
Autor mentem}{%
\lemma{}%
\Bfootnote{%
non-rectum~, %
\bars{doctissimus~\textit{gestr.}} %
autor mentem~\textit{L}
\textbullet~non-rectum autor mentem~\textit{l}
\textbullet~non rectum~, Autor mentem~\textit{E\textsuperscript{1}}}}
suam nondum
\edtext{explicuit, quomodo}{%
\lemma{explicuit}%
\Bfootnote{%
, %
\textit{(1)}~qua nam tunc quoque ambo plana simul totam vim globi %
\textit{(a)}~sustinent, %
\textit{(b)}~sustinerent %
\textit{(2)}~quomodo~\textit{L}
\repeatlemma %
\textit{(1)}~quando~\textit{l} %
\textit{(2)}~quomodo~\textit{Lil}
\repeatlemma, quomodo~\textit{E\textsuperscript{1}}}}
pressiones\indexS{pressio globi} inter ipsa distribuerentur. 
\edlabel{AE_1687_197-198_13a}%
\edtext{}{%
\xxref{AE_1687_197-198_13a}{AE_1687_197-198_13b}%
\lemma{Caeterum \lbrack...\rbrack\ angulos}%
\Cfootnote{%
\indexN{\textls{Vanni}, Giovanni Francesco 1638\textendash1709}%
\textsc{Vanni}, \cite{02000}a.\,a.\,O., I, Prop.~XXI, S.\;29\,f.}}%
Caeterum 
\edtext{ex positis}{%
\lemma{ex}%
\Bfootnote{%
\textit{(1)}~tradit %
\textit{(2)}~positis~\textit{L}}}
\edtext{porro infert,}{%
\lemma{porro}%
\Bfootnote{%
\textit{(1)}~inferre ve
\textit{(2)}~infert~\textit{L}
\repeatlemma infert~\textit{l}
\repeatlemma infert,~\textit{E\textsuperscript{1}}}}
\edtext{etiam velocitates,\indexS{velocitas qua globus in plano descendit} quibus globi\indexS{globus} in plano}{%
\lemma{etiam}%
\Bfootnote{%
\textit{(1)}~velocitatem qua globus in planis
\textit{(2)}~velocitates qua globi in plano~\textit{L}
\repeatlemma velocitates, %
\textit{(1)}~qua globi in plano~\textit{l} %
\textit{(2)}~quibus globi in plano~\textit{Lil}
\repeatlemma velocitates, quibus globi in plano~\textit{E\textsuperscript{1}}}}
inclinato\indexS{planum inclinatum} descendere nituntur,  
\edtext{fore non}{%
\lemma{fore}%
\Bfootnote{%
\textit{(1)}~ut a %
\textit{(2)}~non~\textit{L}}}
reciproce ut longitudines planorum,\indexS{longitudo plani} sed  
\edtext{potius}{%
\lemma{potius}%
\Bfootnote{,~\textit{L}
\repeatlemma~\textit{l}}}
ut eorum
\edtext{angulos;\indexS{angulus elevationis}\edlabel{AE_1687_197-198_13b} \edlabel{AE_1687_197-198_14a}non}{%
\lemma{angulos}%
\Bfootnote{. Non~\textit{l}
\repeatlemma; non~\textit{l}}}%
\edtext{}{%
\xxref{AE_1687_197-198_14a}{AE_1687_197-198_14b}%
\lemma{non vult \lbrack...\rbrack\ descendere}%
\Cfootnote{%
\cite{02000}a.\,a.\,O., Prop.~XXIII, S.\;30\,f.}}
vult  
\edtext{adhiberi globos}{%
\lemma{adhiberi}%
\Bfootnote{%
\textit{(1)}~momenta inaequal %
\textit{(2)}~globos~\textit{L}}}
inaequales\indexS{globi inaequales} in momentis gravium\indexS{momentum gravis} investigandis, quia experientia
\edtext{ostendat,}{%
\lemma{ostendat}%
\Bfootnote{\textit{L~u.~l}
\repeatlemma,~\textit{E\textsuperscript{1}}}}
globos inaequales\indexS{globi inaequales} non aeque velociter descendere,\edlabel{AE_1687_197-198_14b}  quod
\edtext{tamen alii accidentale}{%
\lemma{tamen}%
\Bfootnote{%
alii \bars{tanquam~\textit{gestr.}} accidentale~\textit{L}
\repeatlemma %
\textit{(1)}~allii %
\textit{(2)}~alii~\textit{Lil} %
\href{\Scan{LH_35_09_20_003v}}%
{\lbrack3~v\textsuperscript{o}\rbrack} accidentale~\textit{l}}}
et ab extrinsecis
\edtext{ortum censent.}{%
\lemma{ortum}%
\Bfootnote{%
\textit{(1)}~hic considerare nolunt %
\textit{(2)}~aliorsum %
\textit{(a)}~relegant %
\textit{(b)}~relegantes %
\textit{(3)}~prodientes %
\textit{(4)}~censentes %
\textit{(a)}~aliorsum relegant %
\textit{(b)}~hic considerare nolunt %
\textit{(c)}~aliorsum ablegant.~\textit{L}
\repeatlemma censentes, aliorsum ablegant.~\textit{l}
\repeatlemma censent.~\textit{E\textsuperscript{1}}}}
\pend
\pstart%
\edlabel{AE_1687_197-198_16}%
In \cite{02000}\textit{Exegesi} secunda 
de vecte\indexS{vectis} eo usque progreditur,  
\edlabel{AE_1687_197-198_10a}%
\edtext{}{%
\xxref{AE_1687_197-198_10a}{AE_1687_197-198_10b}%
\lemma{ut neget \lbrack...\rbrack\ commune}%
\Cfootnote{%
\indexN{\textls{Vanni}, Giovanni Francesco 1638\textendash1709}%
\textsc{Vanni}, \cite{02000}a.\,a.\,O., pars II, Prop. XXII, S.\;74\,f.}}%
ut neget  
\edtext{p.~74.}{%
\lemma{}%
\Bfootnote{p.~74. \textit{erg.~L}}}
globorum\indexS{globi qui funiculis tantum connectuntur} qui funiculis\indexS{funiculus}  
\edtext{tantum}{%
\lemma{}%
\Bfootnote{tantum \textit{erg.~L}}}
\edtext{connectuntur, ullum}{%
\lemma{connectuntur}%
\Bfootnote{%
,~%
\textit{(1)}~considerantur ullum debere centrum %
\textit{(2)}~ullum~\textit{L}
\repeatlemma ullum~\textit{l}
\repeatlemma, ullum~\textit{E\textsuperscript{1}}}}
debere considerari centrum gravitatis commune.\indexS{centrum gravitatis commune}\edlabel{AE_1687_197-198_10b} 
Et proinde licet centrum gravitatis\indexS{centrum gravitatis commune}  
\edtext{eorum commune}{%
\lemma{eorum commune}%
\Bfootnote{,~\textit{L}
\repeatlemma~\textit{l}}}
(quod scilicet Geometrae\indexS{Geometrae} hactenus tale appellarunt, punctum scilicet quod
\edtext{rectam, centrum\indexS{centrum gravitatis} singulorum}{%
\lemma{rectam}%
\Bfootnote{%
, centra singulorum~\textit{L}
\repeatlemma centrum %
\textit{(1)}~angulorum~\textit{l} %
\textit{(2)}~singulorum~\textit{Lil}
\repeatlemma, centrum singulorum~\textit{E\textsuperscript{1}}}}
\edtext{conjungentem,}{%
\lemma{conjungentem}%
\Bfootnote{\textit{l}
\repeatlemma,~\textit{L~u.~E\textsuperscript{1}}}}
in ratione  
\edtext{distantiarum corporibus}{%
\lemma{distantiarum}%
\Bfootnote{%
\textit{(1)}~reciproca secat %
\textit{(2)}~corporibus~\textit{L}}}
reciproca secet)
in planis inclinatis\indexS{planum inclinatum} descendere nequeat,\indexS{centrum gravitatis descendere nequiens}
\edlabel{AE_1687_197-198_11a}%
\edtext{}{%
\xxref{AE_1687_197-198_11a}{AE_1687_197-198_11b}%
\lemma{negat \lbrack...\rbrack\ aequilibrio}%
\Cfootnote{%
\cite{02000}a.\,a.\,O., S.\;75.}}%
negat tamen ea esse in
\edtext{aequilibrio,\indexS{aequilibrium}\edlabel{AE_1687_197-198_11b} %
quod hactenus pro indubitato apud staticos\indexS{statici} principio\indexS{principium apud staticos indubitatum} habebatur.}{%
\lemma{aequilibrio}%
\Bfootnote{%
\textit{(1)}~. Caeterum ibidem \bars{denique~\textit{erg.}} pag.~81.
\textit{(a)}~videtur quodammodo in gratiam redire cum communi sententia concedens
\textit{(aa)}~globos duos in diversis planis declivibus positos
\textit{(aaa)}~esse in aequilibrio, seu mutuo se
\textit{(bbb)}~et 
\textit{(aaaa)}~funiculis
\textit{(bbbb)}~funiculo connexos se mutuo sustinere, si
\textit{(bb)}~globis duobus in diversis planis declivibus positis et funiculo connexis ac sese mutuo sustinentibus, ut maneant immoti
\textit{(b)}~ait duos globos 
\textit{(aa)}~funiculis
\textit{(bb)}~funiculo connexos
\textit{(aaa)}~in
\textit{(bbb)}~super diversis planis inclinatis positos, se mutuo sustinere et manere immotos, si sci
\textit{(ccc)}~quorum unus liber pendet 
\textit{(aaaa)}~alter super
\textit{(bbbb)}~alter super plano inclinato est positus, se sustinere, et manere immotos, quando non integer globus in plano inclinato positus, ut
\textit{(aaaaa)}~volunt
\textit{(bbbbb)}~vult communis sententia, sed pars ejus a plano non sustentata, est ut planum ad perpendiculum. Ubi communis sententiae autores dicerent, non partem quidem globi, sed tamen partem momenti ejus
\bars{a plano~\textit{gestr.}}
\textit{(aaaaa-a)}~sustineri
\textit{(bbbbb-b)}~non sustineri. At autor noster partem etiam a trochlea
\textit{(2)}~, quod hactenus pro indubitato apud Staticos principio habebatur.~\textit{L}
\repeatlemma quod hactenus pro indubitato apud Staticos principio habebatur.~\textit{l}
\repeatlemma, \lbrack...\rbrack\ habebatur.~\textit{E\textsuperscript{1}}}}
\edlabel{AE_1687_197-198_12a}%
\edtext{}{%
\xxref{AE_1687_197-198_12a}{AE_1687_197-198_12b}%
\lemma{Quando \lbrack...\rbrack\  sustineri}%
\Cfootnote{%
\cite{02000}a.\,a.\,O., Prop.~XXVIII, S.\;81\,f.}}%
Quando duo  
\edtext{globi\indexS{globus} in planis inclinatis\indexS{planum inclinatum} positi, et}{%
\lemma{globi}%
\Bfootnote{%
\bars{in planis inclinatis positi et~\textit{erg.}}~\textit{L}
\repeatlemma in planis inclinatis positi et~\textit{l}
\repeatlemma in planis inclinatis positi, et~\textit{E\textsuperscript{1}}}}  
\edlabel{AE_1687_197-198_4a}%
\edtext{}{%
\xxref{AE_1687_197-198_4a}{AE_1687_197-198_4b}%
\lemma{funiculo}%
\Bfootnote{%
\textit{(1)}~connecti inter %
\textit{(2)}~connexi super trochlea incedente %
\textit{(a)}~, quorum unus statim %
\textit{(b)}~se mutuo \lbrack...\rbrack\ a trochlea, %
\textit{(aa)}~subsis %
\textit{(bb)}~sustineri %
\textit{(aaa)}~, easque definire %
\textit{(aaaa)}~tentat %
\textit{(bbbb)}~aggreditur. %
\textit{(bbb)}~. \textit{Exegesi} %
\textit{(aaaa)}~3tia de %
\textit{(bbbb)}~tertia~\textit{L}
\repeatlemma connexi, super \bars{trochlea incedente se mutuo sustinent; notat \lbrack...\rbrack\ partem a \textit{erg.~Lil}} trochlea substineri. \textit{Exegesi} tertia~\textit{l}
\repeatlemma \lbrack...\rbrack\ tertia~\textit{E\textsuperscript{1}}}}%
funiculo connexi,\indexS{globi funiculo super trochlea incedente connexi} 
super trochlea\indexS{trochlea} incedente se mutuo sustinent, notat
partem virium\indexS{vis} uniuscujusque globi\indexS{globus} a plano,\indexS{planum inclinatum} partem a trochlea\indexS{trochlea} sustineri.\edlabel{AE_1687_197-198_12b}
\pend
\pstart
\cite{02000}\textit{Exegesi} tertia%
\edlabel{AE_1687_197-198_4b}
de 
\edtext{motu gravium\indexS{motus gravium} \edlabel{AE_1687_197-198_15a}consentit}{%
\lemma{motu}%
\Bfootnote{%
gravium %
\textit{(1)}~concedit %
\textit{(2)}~consentit~\textit{L}}}%
\edtext{}{%
\xxref{AE_1687_197-198_15a}{AE_1687_197-198_15b}%
\lemma{consentit \lbrack...\rbrack\ liberorum}%
\Cfootnote{%
\indexN{\textls{Vanni}, Giovanni Francesco 1638\textendash1709}%
\textsc{Vanni}, \cite{02000}a.\,a.\,O., pars III, S.\;97. Vgl. 
\indexN{\textls{Galilei} (Galilaeus, Galileus), Galileo 1564\textendash1642}%
\textsc{G.~Galilei}, %
\cite{01084}\textit{Discorsi}, %
Giornata Terza, theor.~I\textendash II (\textit{Opere}, Bd.~II, S.\;129\textendash131; \cite{00048}\textit{GO} VIII, S.\;208\textendash213).}}
quidem cum \indexN{\textls{Galilei} (Galilaeus, Galileus), Galileo 1564\textendash1642}%
Galilaeo circa descensum gravium liberorum,\indexS{descensus gravium liberorum}\edlabel{AE_1687_197-198_15b}
\edtext{at non circa descensum in planis inclinatis,\indexS{descensus gravis in plano inclinato}}{%
\lemma{at non \lbrack...\rbrack\ inclinatis}%
\Cfootnote{%
\indexN{\textls{Vanni}, Giovanni Francesco 1638\textendash1709}%
\textsc{Vanni}, \cite{02000}\textit{Exegeses}, pars III \textit{passim}, bes. Prop.~XII-XXII, S.\;108-115.
Vgl. \indexN{\textls{Galilei} (Galilaeus, Galileus), Galileo 1564\textendash1642}%
\textsc{G.~Galilei}, %
\cite{01084}\textit{Discorsi}, %
Giornata Terza, Scholium nach theor.~II bis prop.~XXXVIII (\textit{Opere} Bd.~II, S.\;132\textendash179; \cite{00048}\textit{GO} VIII, S.\;214\textendash267).}}
ut scilicet \cite{02000}\textit{Exegeseos} primae 
\edtext{dogmata tueatur.}{%
\lemma{dogmata}%
\Bfootnote{%
\textit{(1)}~tuentur.~\textit{l} %
\textit{(2)}~tueatur.~\textit{Lil}}}
\pend
\vspace{1em}
\pstart
\noindent
\lbrack%
\textit{Rechnungen in L (Bl.~\href{\Scan{LH_35_09_20_002v}}{2~v\textsuperscript{o}}) ohne erkennbaren Zusammenhang mit dem Text:}\rbrack
\pend
\nobreak
\vspace{0.5em}
\nobreak
\pstart
\begin{edtabularr}
\rule{7.5mm}{0pt}
\edlabel{AE_1687_197-198_2a}%
\edtext{}{%
\xxref{AE_1687_197-198_2a}{AE_1687_197-198_2b}%
\lemma{25 \lbrack...\rbrack\ 48}%
\Cfootnote{%
Die zweite und dritte Spalte lassen sich als Hilfsrechnungen zum Konzept 
\cite{01383}\textit{Leges Concursuum} von Oktober 1686\textendash Februar 1687 
(\textit{LSB}~VIII,~3 N.\;67) 
identifizieren.
Siehe die Datierungsbegründung des Konzepts, \textit{LSB}~VIII,~3 S.\;735 und die Nebenrechnung in der Marginalie zu S.\;742 Z.~2\textendash5.}}%
  18\; & & 25\; &  & 16\; & \hfill \\ 
  4\;\rule[-0.75ex]{0pt}{0pt} & & 6\; &  & 3\; & \hfill \\ 
$\overline{\rule{0pt}{2.4ex}\phantom{22}22}$\; &  & $\overline{\rule{0pt}{2.4ex}\phantom{2}150}$\; &  & $ \overline{\rule{0pt}{2.4ex}\phantom{2}48}$\; & \hfill   \\ 
  9\;\rule[-0.75ex]{0pt}{0pt} &  & 48\; &  & & \hfill  \\ 
  $\overline{\rule{0pt}{2.4ex}\phantom{2}198}$\; &  & $\overline{\rule{0pt}{2.4ex}\phantom{2}198}$\;%
\edlabel{AE_1687_197-198_2b} &  &  & \hfill  
\end{edtabularr}
\pend
\normalsize
